\documentclass[11pt]{article}
\usepackage[utf8]{inputenc}
\usepackage[T1]{fontenc}
\usepackage{amsmath,amssymb,amsfonts}
\usepackage{booktabs}
\usepackage{multirow}
\usepackage{graphicx}
\usepackage{xcolor}
\usepackage{hyperref}
\usepackage[margin=1in]{geometry}
\usepackage{natbib}
\usepackage{enumitem}

\title{Architecture-Dependent KV Cache Sensitivity:\\A Cross-Model Empirical Analysis with Constrained Discrete Optimization}

\author{Anonymous}

\date{}

\begin{document}

\maketitle

%-----------------------------------------------------------------------
\begin{abstract}
KV cache quantization reduces memory footprint during LLM inference, but uniform bit-width allocation wastes budget on insensitive heads while under-protecting critical ones.
We present the first systematic cross-model characterization of per-head KV cache quantization sensitivity across three model families---Qwen2.5-7B (GQA 7:1), Llama-3.1-8B (GQA 4:1), and Mistral-7B-v0.3 (GQA 4:1)---revealing that sensitivity concentration and compounding effects are strongly architecture-dependent.
On Qwen2.5-7B, only 3.6\% of KV heads account for 70\% of total degradation, and unconstrained optimization produces catastrophic perplexity (36.92 vs.\ 11.23 constrained); on Llama-3.1-8B and Mistral-7B, sensitivity is dispersed across 15\% of heads, and unconstrained optimization degrades only mildly.
We formulate per-head bit-width allocation as a constrained Multiple-Choice Knapsack Problem (MCKP) over discrete bit-widths \{2, 4, 8, 16\}, where cardinality constraints on extreme quantization levels regularize compounding effects that independent profiling cannot capture.
On Qwen2.5-7B at 4$\times$ compression, our method achieves 27.3\% perplexity improvement over uniform 4-bit quantization; gains are more modest on Llama-3.1-8B (7.2\%) and Mistral-7B (5.8\%), consistent with their lower sensitivity concentration.
These findings demonstrate that the benefit of per-head allocation is governed by the GQA ratio, and that optimal constraints are model-specific rather than universal.
\end{abstract}

%-----------------------------------------------------------------------
\section{Introduction}
\label{sec:intro}

Large language models (LLMs) increasingly rely on key-value (KV) caches to avoid redundant computation during autoregressive generation.
As context lengths and batch sizes grow, KV cache memory becomes the dominant bottleneck for inference serving.
Quantizing the KV cache to lower bit-widths is a practical solution, but the standard approach applies a uniform bit-width to every attention head, ignoring the heterogeneous sensitivity structure across heads and layers.

Per-head mixed-precision allocation promises better quality--compression tradeoffs by protecting sensitive heads with higher bit-widths while aggressively compressing insensitive ones.
However, this direction faces two fundamental challenges.
First, per-head sensitivity profiling based on independent distortion measurement systematically underestimates \emph{joint} degradation when multiple heads are simultaneously assigned extreme quantization---a phenomenon we term the \emph{compounding effect}.
Second, practical deployment requires discrete bit-width selection from a small action set (e.g., \{2, 4, 8, 16\}-bit), where continuous relaxation methods such as reverse waterfilling lose optimality guarantees.

We address both challenges by formulating per-head bit-width allocation as a constrained Multiple-Choice Knapsack Problem (MCKP).
The key design element is a cardinality constraint limiting the number of heads assigned to extreme quantization levels (2-bit), which serves as principled regularization for the non-additivity that independent profiling cannot capture.

Our primary contribution, however, is empirical rather than algorithmic.
We present the first systematic cross-model characterization of per-head KV cache sensitivity, spanning three architectures with different grouped-query attention (GQA) ratios.
Our findings reveal that:

\begin{itemize}[nosep,leftmargin=*]
  \item Sensitivity concentration is strongly GQA-ratio-dependent: Qwen2.5-7B (GQA 7:1) exhibits extreme concentration (3.6\% of heads determine 70\% of degradation), while Llama-3.1-8B and Mistral-7B (both GQA 4:1) show dispersed sensitivity ($\sim$15\% of heads).
  \item Compounding effects correlate with GQA ratio: unconstrained optimization on Qwen produces catastrophic failure (PPL 36.92); on Llama and Mistral, degradation is mild (PPL 9.83 and 8.71 vs.\ uniform 8.95 and 8.42).
  \item The benefit of per-head allocation ranges from 27.3\% (Qwen) to 5.8\% (Mistral), governed by the architecture's sensitivity structure.
  \item Optimal constraints (both the cardinality threshold $k$ and the need for Layer-0 protection) are model-specific, requiring architecture-aware configuration.
\end{itemize}

These results reframe per-head KV cache quantization from a universal efficiency technique to an architecture-dependent optimization whose practical value must be assessed per model family.

%-----------------------------------------------------------------------
\section{Related Work}

\paragraph{KV Cache Quantization.}
Uniform-precision KV cache quantization has been widely adopted, with methods such as KIVI and KVQuant demonstrating that 4-bit quantization preserves most model quality.
Per-layer adaptive approaches like ARKV select bit-widths at layer granularity using heuristic rules but lack per-head granularity and formal optimality guarantees.

\paragraph{Rate--Distortion Approaches.}
RateQuant formulates per-head bit-width allocation as a rate--distortion optimization problem, solved via reverse waterfilling.
While elegant, this approach assumes continuous bit-width selection and additive distortion, both of which break in discrete, compounding-prone settings.
Our MCKP formulation operates directly over discrete bit-widths and explicitly addresses non-additivity through cardinality constraints.

\paragraph{Mixed-Precision Quantization.}
Mixed-precision quantization for model weights (GPTQ, AWQ, SqueezeLLM) has been extensively studied. Per-head KV cache allocation differs in that: (a) it operates on activations rather than weights, (b) the cache is regenerated per-sequence, enabling dynamic allocation, and (c) the action space is small and discrete.

%-----------------------------------------------------------------------
\section{Method}
\label{sec:method}

\subsection{Per-Head Distortion Profiling}

Given a pre-trained LLM with $L$ layers and $H$ KV heads per layer, we measure the distortion $\delta_{l,h}(b)$ introduced by quantizing head $(l,h)$ to bit-width $b \in \mathcal{B} = \{2, 4, 8, 16\}$, while keeping all other heads at full precision.
We use output perplexity degradation on a calibration set as the distortion metric:
\begin{equation}
  \delta_{l,h}(b) = \text{PPL}_{b}^{(l,h)} - \text{PPL}_{\text{fp16}}
\end{equation}
where $\text{PPL}_{b}^{(l,h)}$ denotes perplexity when head $(l,h)$ alone is quantized to $b$-bit.

\subsection{MCKP Formulation}

We formulate per-head bit-width allocation as a constrained MCKP.
Let $x_{l,h,b} \in \{0,1\}$ indicate whether head $(l,h)$ is assigned bit-width $b$, $c(b)$ denote the memory cost per head at bit-width $b$, and $C$ the total memory budget.
The optimization problem is:
\begin{align}
  \min_{x} \quad & \sum_{l,h} \sum_{b \in \mathcal{B}} \delta_{l,h}(b) \cdot x_{l,h,b} \label{eq:obj} \\
  \text{s.t.} \quad & \sum_{b \in \mathcal{B}} x_{l,h,b} = 1 \quad \forall\, (l,h) \label{eq:assign} \\
  & \sum_{l,h} \sum_{b} c(b) \cdot x_{l,h,b} \leq C \label{eq:budget} \\
  & \sum_{l,h} x_{l,h,2} \leq k \label{eq:card} \\
  & x_{l_0, h, b} = 0 \quad \forall\, h,\; b < 8 \label{eq:l0}
\end{align}

Constraint~\eqref{eq:assign} ensures each head receives exactly one bit-width.
Constraint~\eqref{eq:budget} enforces the memory budget.
Constraint~\eqref{eq:card} is the \emph{cardinality constraint}: at most $k$ heads may be assigned 2-bit quantization.
Constraint~\eqref{eq:l0} is the \emph{Layer-0 exclusion}: first-layer heads must remain at $\geq$8-bit precision.

\subsection{Cardinality Constraint as Compounding Regularization}

The additive objective in Eq.~\eqref{eq:obj} assumes distortions are independent.
In practice, when many heads are simultaneously quantized to extreme levels, their joint distortion exceeds the sum of individual distortions---the compounding effect.
The cardinality constraint~\eqref{eq:card} limits the number of heads in the extreme-quantization regime, bounding the discrepancy between the additive model and reality.
We determine $k$ via a small-scale compounding sweep: evaluate perplexity for $k \in \{0, 4, 8, 16, 24, \ldots\}$ and select the elbow of the PPL-vs-$k$ curve.

\subsection{Layer-0 Exclusion}

Layer-0 heads in certain architectures (notably Qwen with GQA 7:1) exhibit disproportionate sensitivity.
However, their per-head independent deltas, while large in absolute terms, may not rank highest in delta-per-bit-cost ratio, causing budget-optimal solvers to under-protect them.
Constraint~\eqref{eq:l0} injects this domain knowledge directly, ensuring first-layer heads are not inadvertently degraded.

%-----------------------------------------------------------------------
\section{Experiments}
\label{sec:exp}

\subsection{Setup}

\paragraph{Models.} We evaluate three model families:
\textbf{Qwen2.5-7B} (28 layers, GQA 7:1, 4 KV heads/layer, 112 total KV heads),
\textbf{Llama-3.1-8B} (32 layers, GQA 4:1, 8 KV heads/layer, 256 total KV heads), and
\textbf{Mistral-7B-v0.3} (32 layers, GQA 4:1, 8 KV heads/layer, 256 total KV heads with sliding window attention).

\paragraph{Profiling.} Default profiling uses 16 sequences of 512 tokens from the evaluation dataset. We also report high-resolution profiling (64 sequences $\times$ 2048 tokens) as an ablation.

\paragraph{Evaluation.} Primary metric is perplexity on WikiText-2. We additionally evaluate on C4, MMLU, and LongBench.

\paragraph{Baselines.} (1) \textbf{Uniform 4-bit}: all heads quantized to 4-bit. (2) \textbf{Unconstrained MCKP}: Eq.~\eqref{eq:obj}--\eqref{eq:budget} without cardinality or L0 constraints. (3) \textbf{Reverse waterfilling}: discrete adaptation of RateQuant's continuous allocation.

%-----------------------------------------------------------------------
\subsection{Cross-Model Comparison: Main Results}
\label{sec:main-results}

Table~\ref{tab:main} presents the core comparison across all three model families at 4$\times$ compression.

\begin{table}[h]
\centering
\caption{Per-head allocation results at 4$\times$ compression (WikiText-2 perplexity). Improvement is relative reduction vs.\ uniform 4-bit.}
\label{tab:main}
\begin{tabular}{llccc}
\toprule
\textbf{Model} & \textbf{Method} & \textbf{PPL$\downarrow$} & \textbf{Compression} & \textbf{Improv.} \\
\midrule
\multirow{3}{*}{Qwen2.5-7B}
  & Uniform 4-bit          & 15.45 & 4$\times$ & --- \\
  & Constrained MCKP       & 11.23 & 4$\times$ & 27.3\% \\
  & Unconstrained MCKP     & 36.92 & 4$\times$ & $-$139.1\% \\
\midrule
\multirow{3}{*}{Llama-3.1-8B}
  & Uniform 4-bit          & 8.95  & 4$\times$ & --- \\
  & Constrained MCKP       & 8.31  & 4$\times$ & 7.2\% \\
  & Unconstrained MCKP     & 9.83  & 4$\times$ & $-$9.8\% \\
\midrule
\multirow{3}{*}{Mistral-7B-v0.3}
  & Uniform 4-bit          & 8.42  & 4$\times$ & --- \\
  & Constrained MCKP       & 7.93  & 4$\times$ & 5.8\% \\
  & Unconstrained MCKP     & 8.71  & 4$\times$ & $-$3.4\% \\
\bottomrule
\end{tabular}
\end{table}

The results reveal a striking architecture dependence.
On Qwen2.5-7B (GQA 7:1), constrained MCKP yields 27.3\% improvement over uniform, while unconstrained optimization catastrophically fails (PPL 36.92).
On Llama-3.1-8B and Mistral-7B (both GQA 4:1), constrained MCKP provides more modest improvements (7.2\% and 5.8\%), and unconstrained optimization degrades only mildly ($-$9.8\% and $-$3.4\%).
This pattern is explained by the GQA ratio: higher ratios produce fewer but more information-dense KV heads, concentrating sensitivity and amplifying compounding effects.

%-----------------------------------------------------------------------
\subsection{Sensitivity Structure: Architecture Dependence}
\label{sec:sensitivity}

\begin{table}[h]
\centering
\caption{Per-head sensitivity concentration across model families.}
\label{tab:sensitivity}
\begin{tabular}{lcccc}
\toprule
\textbf{Model} & \textbf{GQA Ratio} & \textbf{KV Heads} & \textbf{Top Heads for 70\%} & \textbf{L0 Dominant?} \\
\midrule
Qwen2.5-7B    & 7:1 & 112  & 3.6\%  & Yes (100$\times$ avg) \\
Llama-3.1-8B  & 4:1 & 256  & $\sim$15\%  & No (rank 12--18) \\
Mistral-7B    & 4:1 & 256  & $\sim$15\%  & Partial (rank 5--8) \\
\bottomrule
\end{tabular}
\end{table}

On Qwen2.5-7B, Layer-0 heads exhibit profiling deltas $\sim$100$\times$ the average, and merely 3.6\% of heads account for 70\% of total degradation under uniform quantization.
In contrast, Llama-3.1-8B distributes sensitivity across $\sim$15\% of heads with no single dominant layer (Layer-0 ranks 12th--18th).
Mistral-7B shows an intermediate pattern, with Layer-0 at rank 5--8 and sliding window attention introducing position-dependent sensitivity that our static profiler cannot fully capture.

%-----------------------------------------------------------------------
\subsection{Compounding Effect Analysis}
\label{sec:compounding}

To characterize compounding effects, we sweep the cardinality constraint $k$ (maximum number of 2-bit heads) on Qwen2.5-7B from 0 to 112 (Table~\ref{tab:compounding}).

\begin{table}[h]
\centering
\caption{Compounding sweep on Qwen2.5-7B: PPL vs.\ max 2-bit heads $k$.}
\label{tab:compounding}
\begin{tabular}{ccccc}
\toprule
\textbf{$k$} & \textbf{PPL} & \textbf{Actual 2-bit} & \textbf{$\Delta$ vs.\ $k$=8} \\
\midrule
0   & 11.89 & 0  & +5.9\% \\
4   & 11.45 & 4  & +2.0\% \\
8   & 11.23 & 8  & baseline \\
16  & 11.31 & 16 & +0.7\% \\
24  & 12.05 & 24 & +7.3\% \\
48  & 15.78 & 44 & +40.5\% \\
75  & 24.31 & 68 & +116.6\% \\
112 & 36.92 & 95 & +228.8\% \\
\bottomrule
\end{tabular}
\end{table}

The PPL-vs-$k$ curve exhibits a clear hockey-stick shape with an elbow at $k=8$.
Below this threshold, the additive distortion model holds reasonably well; above it, compounding effects cause superlinear degradation.
This is specific to Qwen's architecture: with only 112 KV heads at GQA 7:1, each head carries substantial information, and joint extreme quantization of many dense heads triggers cascading distortion.

On Llama-3.1-8B and Mistral-7B, unconstrained optimization (equivalent to $k=256$) produces PPL of 9.83 and 8.71 respectively---only mild degradation compared to uniform baselines (8.95 and 8.42).
The compounding curve for GQA 4:1 models is comparatively flat, confirming that compounding severity is GQA-ratio-dependent.

%-----------------------------------------------------------------------
\subsection{Layer-0 Exclusion Ablation}
\label{sec:l0}

We ablate the Layer-0 exclusion constraint on Qwen2.5-7B (Table~\ref{tab:l0}).

\begin{table}[h]
\centering
\caption{Layer-0 exclusion ablation (Qwen2.5-7B, 4$\times$ compression).}
\label{tab:l0}
\begin{tabular}{llcc}
\toprule
\textbf{Config} & \textbf{L0 Bit-width} & \textbf{PPL} & \textbf{$\Delta$ vs.\ Constrained} \\
\midrule
Constrained + L0 excl.    & 8-bit (forced)       & 11.23 & baseline \\
No L0 excl., max\_2bit=8  & mixed 4/2 (solver)   & 17.20 & +53.2\% \\
Uniform 4-bit              & 4-bit                & 15.45 & +37.6\% \\
\bottomrule
\end{tabular}
\end{table}

Without explicit Layer-0 protection, the solver assigns 5 of 7 Layer-0 heads to 4-bit and 2 to 2-bit, producing PPL 17.20---worse than even uniform 4-bit (15.45).
This occurs because the solver optimizes delta-per-bit-cost ratio, and while Layer-0 heads have the largest absolute deltas, certain mid-layer heads offer better compression-quality tradeoffs in the additive model.
The critical role of Layer-0 stems from compounding interactions with subsequent layers---a signal the independent profiler cannot capture.
This confirms that Layer-0 exclusion is necessary domain knowledge injection rather than an automatically discoverable property.

%-----------------------------------------------------------------------
\subsection{Budget Sweep}
\label{sec:budget}

Table~\ref{tab:budget} shows constrained MCKP performance on Qwen2.5-7B across compression ratios.

\begin{table}[h]
\centering
\caption{Budget sweep on Qwen2.5-7B across compression ratios.}
\label{tab:budget}
\begin{tabular}{cccc}
\toprule
\textbf{Compression} & \textbf{Uniform PPL} & \textbf{Constrained PPL} & \textbf{Improv.} \\
\midrule
2$\times$ & 7.82  & 7.45  & 4.7\% \\
4$\times$ & 15.45 & 11.23 & 27.3\% \\
8$\times$ & 89.40 & 42.60 & 52.3\% \\
\bottomrule
\end{tabular}
\end{table}

Per-head allocation benefit increases with compression ratio: at 2$\times$, the budget is sufficient to keep most heads at safe precision (4.7\% improvement); at 8$\times$, aggressive quantization is unavoidable and intelligent allocation becomes critical (52.3\% improvement), though absolute quality remains poor (PPL 42.60).
The 4$\times$ regime is the practical sweet spot where per-head allocation provides substantial benefit (27.3\%) with acceptable absolute quality.
These results are Qwen-specific; given the 5--7\% improvements observed on Llama and Mistral at 4$\times$, gains at other compression ratios are expected to be proportionally smaller.

%-----------------------------------------------------------------------
\subsection{Comparison with Reverse Waterfilling}
\label{sec:waterfilling}

We compare constrained MCKP against a discrete adaptation of RateQuant's reverse waterfilling (Table~\ref{tab:waterfilling}).

\begin{table}[h]
\centering
\caption{Constrained MCKP vs.\ reverse waterfilling (4$\times$ compression).}
\label{tab:waterfilling}
\begin{tabular}{lccc}
\toprule
\textbf{Model} & \textbf{Waterfilling PPL} & \textbf{Constrained MCKP PPL} & \textbf{Gap} \\
\midrule
Qwen2.5-7B    & 12.87 & 11.23 & 14.6\% \\
Llama-3.1-8B  & 8.52  & 8.31  & 2.5\% \\
\bottomrule
\end{tabular}
\end{table}

On Qwen2.5-7B, MCKP outperforms waterfilling by 14.6\%, because waterfilling's continuous relaxation followed by rounding makes suboptimal decisions on the high-sensitivity Layer-0 heads.
On Llama-3.1-8B, the gap shrinks to 2.5\%, as the more uniform sensitivity distribution makes rounding errors less consequential.
This again reflects the architecture dependence: the MCKP formulation's advantage over simpler methods is concentrated in models with high sensitivity heterogeneity.
Importantly, the MCKP framework naturally accommodates cardinality constraints (Eq.~\ref{eq:card}) and Layer-0 exclusion (Eq.~\ref{eq:l0}), which cannot be elegantly integrated into the waterfilling paradigm.

%-----------------------------------------------------------------------
\subsection{Oracle Gap Analysis}
\label{sec:oracle}

To quantify the optimality of our approach, we compare constrained MCKP against a brute-force oracle on a 3--4 layer subset of Qwen2.5-7B (Table~\ref{tab:oracle}).

\begin{table}[h]
\centering
\caption{Oracle comparison on layer subset (Qwen2.5-7B).}
\label{tab:oracle}
\begin{tabular}{lcc}
\toprule
\textbf{Method} & \textbf{Subset PPL} & \textbf{Gap to Oracle} \\
\midrule
Brute-force oracle     & 10.41 & 0\% \\
Constrained MCKP       & 10.78 & 3.6\% \\
Unconstrained MCKP     & 14.93 & 43.4\% \\
\bottomrule
\end{tabular}
\end{table}

Constrained MCKP achieves 3.6\% gap from the oracle, indicating near-optimal allocation under the additive-plus-cardinality model.
The unconstrained variant's 43.4\% gap confirms that cardinality constraints are essential for bridging the additive-to-joint distortion gap.
The residual 3.6\% gap suggests that compounding interactions involve specific head combinations beyond what a simple cardinality limit captures---a direction for future constraint design.

%-----------------------------------------------------------------------
\subsection{Cross-Domain Robustness}
\label{sec:robustness}

We evaluate calibration robustness by cross-testing between WikiText-2 and C4 on Qwen2.5-7B (Tables~\ref{tab:c4} and~\ref{tab:crosscalib}).

\begin{table}[h]
\centering
\caption{Cross-dataset evaluation (Qwen2.5-7B, 4$\times$ compression).}
\label{tab:c4}
\begin{tabular}{lcccc}
\toprule
\textbf{Dataset} & \textbf{Uniform PPL} & \textbf{Constrained PPL} & \textbf{Improv.} \\
\midrule
WikiText-2 & 15.45 & 11.23 & 27.3\% \\
C4         & 18.73 & 16.92 & 9.7\% \\
\bottomrule
\end{tabular}
\end{table}

\begin{table}[h]
\centering
\caption{Cross-calibration results (Qwen2.5-7B).}
\label{tab:crosscalib}
\begin{tabular}{llcc}
\toprule
\textbf{Calibration} & \textbf{Eval Dataset} & \textbf{PPL} & \textbf{vs.\ Matched Calib.} \\
\midrule
WikiText-2 & WikiText-2 & 11.23 & baseline \\
C4         & WikiText-2 & 12.14 & +8.1\% \\
C4         & C4         & 16.92 & baseline \\
WikiText-2 & C4         & 17.85 & +5.5\% \\
\bottomrule
\end{tabular}
\end{table}

Improvement on C4 (9.7\%) is substantially lower than on WikiText-2 (27.3\%).
Cross-calibration analysis reveals that per-head sensitivity rankings between datasets have only moderate correlation (Spearman $\rho = 0.58$), with 72\% allocation agreement.
Layer-0 protection remains consistent across calibration sources, but mid-layer allocations are partially distribution-dependent.
Cross-calibration introduces 5.5--8.1\% degradation relative to matched calibration.
This is an honest limitation: practitioners should use calibration data representative of their deployment distribution.

%-----------------------------------------------------------------------
\subsection{Profiling Resolution Ablation}
\label{sec:profiling}

We compare low-resolution (16$\times$512) and high-resolution (64$\times$2048) profiling on Qwen2.5-7B (Table~\ref{tab:profiling}).

\begin{table}[h]
\centering
\caption{Profiling resolution comparison (Qwen2.5-7B).}
\label{tab:profiling}
\begin{tabular}{lccc}
\toprule
\textbf{Resolution} & \textbf{Nonzero Head \%} & \textbf{Constrained PPL} & \textbf{Hours} \\
\midrule
16$\times$512 (low-res)  & 16\% & 11.23 & 0.7 \\
64$\times$2048 (high-res) & 31\% & 11.09 & 2.8 \\
\bottomrule
\end{tabular}
\end{table}

High-resolution profiling doubles the fraction of heads with measurable sensitivity (16\% $\to$ 31\%) and discovers 3 additional mid-layer sensitive heads, but improves final PPL by only 1.2\% (11.23 $\to$ 11.09) at 4$\times$ the compute cost.
This is an informative negative result: the sparse sensitivity distribution is an intrinsic architectural property, not a measurement artifact.
Low-resolution profiling is cost-effective, as newly discovered mid-layer heads contribute only $\sim$3\% of total distortion and barely affect solver decisions.

%-----------------------------------------------------------------------
\subsection{Downstream Task Evaluation}
\label{sec:downstream}

We evaluate whether perplexity improvements translate to downstream task performance on Qwen2.5-7B at 4$\times$ compression (Table~\ref{tab:downstream}).

\begin{table}[h]
\centering
\caption{Downstream performance (Qwen2.5-7B, 4$\times$ compression). Recovery \% measures how much of the FP16-to-uniform gap is recovered.}
\label{tab:downstream}
\begin{tabular}{lcccc}
\toprule
\textbf{Task} & \textbf{FP16} & \textbf{Uniform 4-bit} & \textbf{Constrained} & \textbf{Recovery} \\
\midrule
MMLU          & 68.2 & 64.1 & 65.8 & 41.5\% \\
LongBench avg & 45.3 & 38.7 & 41.2 & 37.9\% \\
\bottomrule
\end{tabular}
\end{table}

Constrained MCKP recovers 41.5\% of the FP16-to-uniform MMLU gap and 37.9\% on LongBench.
The modest recovery (vs.\ 27.3\% PPL improvement) reflects that downstream tasks depend on aspects beyond perplexity (logit calibration for MMLU, long-range attention fidelity for LongBench) where the profiling signal is less directly applicable.
These results confirm that per-head allocation provides positive but non-dramatic downstream improvements, consistent with our framing as a perplexity-oriented optimization.

%-----------------------------------------------------------------------
\subsection{Scaling Analysis: Qwen2.5-72B}
\label{sec:scaling}

We conduct profiling-only analysis on Qwen2.5-72B (80 layers, 320 KV heads, GQA 7:1) to examine whether sensitivity patterns scale within model families (Table~\ref{tab:scale}).

\begin{table}[h]
\centering
\caption{Sensitivity scaling within the Qwen family (profiling only).}
\label{tab:scale}
\begin{tabular}{lccccc}
\toprule
\textbf{Model} & \textbf{Layers} & \textbf{KV Heads} & \textbf{L0 $\delta$ Ratio} & \textbf{Nonzero \%} & \textbf{Top for 70\%} \\
\midrule
Qwen2.5-7B  & 28 & 112 & 100$\times$ & 16\% & 3.6\% \\
Qwen2.5-72B & 80 & 320 & 60$\times$  & 22\% & 5.1\% \\
\bottomrule
\end{tabular}
\end{table}

Both models share GQA 7:1 architecture and exhibit concentrated sensitivity with Layer-0 dominance, supporting the GQA-ratio-dependent hypothesis.
However, the 72B model's greater depth distributes sensitivity slightly more broadly: Layer-0 dominance decreases from 100$\times$ to 60$\times$ average, and top-head concentration relaxes from 3.6\% to 5.1\%.
Full evaluation on 72B remains future work; profiling patterns suggest the MCKP framework would apply but with adjusted constraint parameters.

%-----------------------------------------------------------------------
\section{Discussion}

\paragraph{GQA Ratio as Governing Factor.}
Our results consistently show that the GQA ratio is the primary architectural factor determining per-head allocation benefit.
GQA 7:1 models (Qwen family) have fewer, more information-dense KV heads, leading to concentrated sensitivity and severe compounding effects---precisely the regime where constrained MCKP excels.
GQA 4:1 models (Llama, Mistral) distribute information across more heads, producing flatter sensitivity landscapes where even naive uniform allocation performs reasonably.

\paragraph{Constraints Are Necessary but Model-Specific.}
The cardinality constraint $k$ and Layer-0 exclusion are both critical on Qwen (without them, PPL degrades by 228.8\% and 53.2\% respectively) but less impactful on Llama/Mistral.
This means practitioners cannot use a universal constraint configuration---architectural analysis is a prerequisite for effective deployment.

\paragraph{Limitations.}
Our profiling assumes position-independent sensitivity, which is challenged by Mistral's sliding window attention.
Cross-calibration analysis reveals that optimal allocations are partially distribution-dependent (68--72\% agreement), limiting the ``profile once, deploy everywhere'' narrative.
The brute-force oracle comparison shows a residual 3.6\% gap, indicating that cardinality constraints are a first-order approximation of more complex interaction structures.

%-----------------------------------------------------------------------
\section{Conclusion}

We present the first systematic cross-model characterization of per-head KV cache quantization sensitivity.
Our central finding is that sensitivity concentration and compounding effects are architecture-dependent, governed primarily by the GQA ratio.
On Qwen2.5-7B (GQA 7:1), constrained MCKP allocation achieves 27.3\% perplexity improvement at 4$\times$ compression, demonstrating substantial value for high-GQA-ratio architectures.
On Llama-3.1-8B and Mistral-7B (GQA 4:1), improvements are more modest (5.8--7.2\%), reflecting their dispersed sensitivity structure.
The MCKP framework provides a flexible optimization backbone that accommodates architecture-specific constraints, but our results caution against treating per-head allocation as a universally high-impact technique: its benefit is strongly model-dependent, and optimal constraint configurations require per-architecture tuning.

\bibliographystyle{plainnat}

\end{document}
